\section{Discovery Saturation Across the Phenome}\label{sec:results}

\begin{figure}[t!]
    \centering
    {{fig:discovery_cycle}}
    \caption{The genetic discovery cycle across the phenome. (a) Schematic illustrating the expected discovery metric and saturation detection approach. (b-e) Exemplar traits representing different discovery stages: (b) Parkinson's disease (early, power-limited---all curves linear), (c) Asthma (mid-stage, annotation-limited), (d) Type 2 diabetes (saturated, pathway-concentrated---indirect diverges from direct), (e) Height (saturated, distributed---all curves saturate together). (f) Distribution of discovery status across {{NUM_TRAITS_SATURATION_ANALYSIS}} common traits.}
    \label{fig:discovery-cycle}
\end{figure}

Across {{NUM_TRAITS_SATURATION_ANALYSIS}} common traits, {{PERCENT_TRAITS_SATURATED}}\% showed evidence of GWAS saturation for direct genetic support (Figure~\ref{fig:discovery-cycle}f). The remaining traits split between power-limited ({{NUM_TRAITS_POWER_LIMITED}} traits) and annotation-limited ({{NUM_TRAITS_ANNOTATION_LIMITED}} traits). Among saturated traits, we observed two distinct architectural patterns based on the relationship between direct and indirect discovery trajectories.

% Exemplar 1: Power-limited
Parkinson's disease exemplifies a trait in the early discovery stage (Figure~\ref{fig:discovery-cycle}b). All discovery curves---direct, indirect, combined, and the number of independent loci---increase linearly with sample size, with no evidence of saturation. The parallel trajectories indicate that both disease-associated and disease-relevant genes are being discovered at comparable rates, and continued GWAS investment will yield new loci that in turn enable identification of additional relevant genes through indirect support.

% Exemplar 2: Annotation-limited
Asthma represents a mid-stage trait where direct discovery is progressing but indirect support lags (Figure~\ref{fig:discovery-cycle}c). Direct support shows early signs of curvature while the number of independent loci continues to grow linearly. Critically, indirect support accumulates more slowly than direct support, suggesting that current functional annotations incompletely capture asthma-relevant biology. For such traits, investment in disease-relevant annotations---airway epithelial expression programs or immune pathway curation---may yield greater returns than additional GWAS sample size alone.

% Exemplar 3: Saturated, pathway-concentrated
Type 2 diabetes exemplifies a saturated trait with pathway-concentrated genetic architecture (Figure~\ref{fig:discovery-cycle}d). Both direct and indirect support show clear exponential saturation, reaching their asymptotes by approximately 40\% of the full sample size. However, combined support diverges upward from direct support, and the number of independent loci continues increasing linearly even as per-gene discovery saturates. This pattern indicates that T2D biology is organized around discrete pathways: once anchored by associated genes, these pathways propagate indirect support to functionally related genes that lack direct associations. For T2D, the discovery frontier lies in understanding pathway relationships rather than finding additional GWAS loci.

% Exemplar 4: Saturated, distributed
Height presents a contrasting saturation pattern (Figure~\ref{fig:discovery-cycle}e). Like T2D, both direct and indirect support saturate early---by approximately 30\% of the full sample size. Unlike T2D, all curves saturate together without divergence: direct, indirect, and combined support reach similar asymptotes. This pattern reflects height's highly polygenic, distributed genetic architecture \cite{yengo_saturated_2022}: thousands of genes contribute small effects through diverse biological processes, and shared functional annotations do not strongly connect associated genes to additional relevant genes beyond those already detected.

% ANALYSIS IDEA: Temporal validation - traits classified as "power-limited" in older GWAS
% should show more new discoveries in newer GWAS than traits classified as "saturated."

% FIGURE SUGGESTION: Supplementary figure showing all ~1000 traits in 2D space
% (saturation degree vs. direct/indirect divergence), colored by disease category.
